\chapter{Introduction}\label{introduction}

\section{Motivation}\label{motivation}

The motivation for Intelli-Trading Fours (IT4s) takes inspiration from
an improvisational practice in Jazz known as ``trading-fours'' where
musicians alternate short improvised phrases in an exchange with each
other. This format enforces a cycle in which listening, interpretation
and response are all derivative of each other, introducing an element of
conversation to musical expression.

This interaction model demands close collaboration and near-instant
responsive creativity as performers. Its success is predicated on the
performers' mutual ability to contextually interpret incoming musical
material and generate a musically coherent reply within a constrained
time window. Interdisciplinary research suggests that interactive
improvisation engages distinct cognitive and neural processes, affirming
this project's emphasis on mutual responsiveness and temporal alignment
in music dialogue \citep{donnay2014neural}.

From a computational perspective, IT4s provides a framework for human-AI
interaction that is both practical and structured. Performances are
segmented into explicit turns comprising input capture, analysis,
decision-making and response generation. This allows the system to
distinguish and interpret each phase through specialised, extensible
subsystems. This separation also reduces the ambiguity associated with
continuous co-improvisation \citep{assayag2006skini}, instead
providing a closed-loop interaction architecture where the system
repeatedly listens, processes, responds and re-enters a waiting state.

This model forms the foundation of the system design: the implementation
formalises interaction as a sequence of alternating human and AI turns,
positioning the system as an interpretable collaborator operating in
real time rather than as a passive offline generator. Rather than
treating generation as an isolated musical artefact, IT4s situates each
response within a repeated temporal exchange, reflecting accounts of
computer music that emphasise musical structure as emerging across
interacting temporal levels \citep{vaggione1993determinism} and interactive systems
that foreground performer-conditioned response (Pachet, 2002; Pasquier
et al., 2016).

This work therefore investigates not only how an AI system can generate
rhythmic material, but more pertinently, how it can meaningfully
participate within a structured, time-sensitive interaction loop that
mirrors human musical dialogue.

\section{Research Problem}\label{research-problem}

Designing an interactive musical system capable of participating in
real-time rhythmic improvisation presents a multi-stage technical
challenge. Real-time systems must sustain a state of temporal coupling
with performer input. Unlike offline or non-interactive generative
approaches, these systems must: capture live musical events, derive an
interpretable representation from them, generate response decisions and
produce the audible output before returning to a ready state quickly
enough for the exchange to remain musically coherent.

Raw performance is not immediately suitable for musical reasoning within
an Human Computer Interaction (HCI) like it might be between two human
performers; this forms a central difficulty for the system. Hit events
or audio signals need transforming into a symbolic representation that
preserves musical properties such as timing, intensity and phrase
boundaries whilst remaining tractable for computation. Research into
both computational music analysis and rhythm perception show rhythmic
interpretation to be more than a binary matter of counting events;
properties such timing, salience, meter and grouping all contribute to
the way rhythmic material is construed \citep{longuetHiggins1984rhythmic,temperley2009music,levitin2018psychology}.

This aforementioned representation subsequently informs response
decisions. In the case of Intelli-Trading Fours, this necessitates an
interpretable planning stage where response behaviour is variably
influenced by analysed features in a controllable manner. Within this
approach lies emphasis on ensuring the relationship between human input
and system response observable and musically justifiable; thus
distinguishing this approach from others that solely prioritise
autonomous or opaque generation when considering the production of
rhythmic material \citep{pachet2002continuator,pasquier2016introduction,smith2022humanai}.

Finally, this interaction must be sustained across repeated cycles.
After musically responding, the system needs to be poised for the next
human turn without accumulating stale state or losing synchronisation.
Therefore positioning this project as a response to the following
problem:

\emph{\textbf{How can a closed-loop, real-time system capture and
represent human rhythmic input, analyse it, plan an interpretable
response, realise an audible output and sustain this process across
repeated turn-taking cycles?}}

Intelli-Trading Fours addresses this problem through the implementation
of a closed-loop rhythmic interaction pipeline spanning Bela input
capture across OSC with Unity-based turn orchestration, symbolic
rhythmic analysis, response planning, interim response realisation,
ChucK/Chunity playback and loop closure.

\section{Aims and Research Questions}\label{aims-and-research-questions}

The aim of this report is to design, implement and evaluate a
closed-loop prototype for turn-based human-AI rhythmic interaction.
Intelli-Trading Fours focuses on the real-time capture of human rhythmic
input, its conversion into a symbolic representation, the extraction of
interpretable rhythmic features, response planning, audible playback and
loop closure across repeated interaction cycles. The project does not
aim to deliver a fully completed autonomous drumming agent; rather, it
investigates the technical architecture required for a system to
participate meaningfully in a structured rhythmic exchange.

To support this aim, the project has the following objectives:

\begin{itemize}
\item
  To implement a real-time input pipeline capable of capturing human
  rhythmic performance through embedded hardware and OSC communication.
\item
  To transform captured hit events into a symbolic rhythmic
  representation suitable for analysis, planning and playback.
\item
  To design rhythmic analysis modules that extract musically relevant
  features such as density, energy, salient anchors, end activity and
  segment-level activity profiles.
\item
  To implement an interpretable response-planning stage in which
  analysed rhythmic features influence response behaviour in a
  controllable and observable manner.
\item
  To integrate interim response realisation and ChucK/Chunity playback
  within a closed-loop turn-taking architecture that can return to a
  ready state after an AI response.
\item
  To evaluate the implemented system through tests, diagnostic tools and
  Unity runtime observability, while identifying the limitations of the
  current response-generation layer.
\end{itemize}

The dissertation is therefore guided by the following research
questions:

\begin{itemize}
\item
  RQ1: How can a real-time system structure repeated turn-based rhythmic
  interaction between a human drummer and an AI agent?
\item
  RQ2: How can human rhythmic input be represented and analysed in a
  form suitable for response planning?
\item
  RQ3: How can interpretable planning connect analysed rhythmic features
  to musically justifiable response behaviour?
\item
  RQ4: What are the technical limitations of closing the interaction
  loop using an interim transformation-based response layer?
\end{itemize}

These questions reflect the implemented scope of Intelli-Trading Fours.
They focus the dissertation on real-time interaction architecture,
rhythmic representation, analysis, planning and loop closure, while
leaving the completion of a richer generative response-realisation
pipeline as a defined area for future work.

\section{Implemented Scope and Deferred Scope}\label{implemented-scope-and-deferred-scope}

Intelli-Trading Fours is implemented as a closed-loop prototype for
turn-based human-AI rhythmic interaction whose system architecture
supports its titular full interaction cycle. The exchange required for a
Trading Fours style cycle: human input capture, symbolic rhythmic
representation compilation, analysis that informs response planning,
audible AI response played back and finally a return to a state ready
for human input capture again. This establishes the project as an
interactive real-time system rather than a collection of disconnected
generation components.

The core infrastructure necessary to support this interaction loop is
enabled through the following layers:

\begin{itemize}
\item
  Input
\item
  Representation
\item
  Analysis \& Planning
\item
  Output
\item
  Observability
\end{itemize}

These infrastructural components form the technical foundation of the
system and are the primary focus of the dissertation.

The response-generation pipeline is deliberately limited in scope.
Whilst the system implements an interpretable response-planning layer
and a prototype structural generation layer with planning logic and
skeleton-oriented response design, the final response realisation stage,
in which planned intent and structural skeletons are interpreted and
converted into a fully developed rhythmic response, remains incomplete.
In the current scope, audible AI responses are produced using an interim
transformation-based mechanism. This design prioritises evaluable
decision-making over expressive synthesis, leaving richer generative
response realisation as a defined area for future work.

This distinction is central to this report. The project should be
understood as a closed-loop rhythmic interaction prototype whose main
contribution lies in the architecture, representation, analysis,
planning, playback integration and observability required to support
repeated human-AI rhythmic exchange. The partial state of the final
generation layer is treated as an implementation limitation and future
development path, rather than as a central measure of the system's
success.

A scope summary is shown in Table~\ref{tab:scope-summary}.

\begin{longtable}[]{@{}
  >{\raggedright\arraybackslash}p{(\linewidth - 4\tabcolsep) * \real{0.4433}}
  >{\raggedright\arraybackslash}p{(\linewidth - 4\tabcolsep) * \real{0.1799}}
  >{\raggedright\arraybackslash}p{(\linewidth - 4\tabcolsep) * \real{0.3597}}@{}}
\caption{Project scope summary}\label{tab:scope-summary}\\
\toprule\noalign{}
\begin{minipage}[b]{\linewidth}\raggedright
\textbf{Area}
\end{minipage} & \begin{minipage}[b]{\linewidth}\raggedright
\textbf{Status}
\end{minipage} & \begin{minipage}[b]{\linewidth}\raggedright
\textbf{Role in Dissertation}
\end{minipage} \\
\begin{minipage}[b]{\linewidth}\raggedright
Skeleton-oriented structural generation
\end{minipage} & \begin{minipage}[b]{\linewidth}\raggedright
Prototype
\end{minipage} & \begin{minipage}[b]{\linewidth}\raggedright
Prototype architecture
\end{minipage} \\
\begin{minipage}[b]{\linewidth}\raggedright
Audible response generation
\end{minipage} & \begin{minipage}[b]{\linewidth}\raggedright
Interim implementation
\end{minipage} & \begin{minipage}[b]{\linewidth}\raggedright
Limitation-aware implementation
\end{minipage} \\
\begin{minipage}[b]{\linewidth}\raggedright
Turn-taking orchestration
\end{minipage} & \begin{minipage}[b]{\linewidth}\raggedright
Implemented
\end{minipage} & \begin{minipage}[b]{\linewidth}\raggedright
Core system architecture
\end{minipage} \\
\begin{minipage}[b]{\linewidth}\raggedright
Bela/OSC input capture + Hit buffering and turn capture
\end{minipage} & \begin{minipage}[b]{\linewidth}\raggedright
Implemented
\end{minipage} & \begin{minipage}[b]{\linewidth}\raggedright
Implementation chapter
\end{minipage} \\
\begin{minipage}[b]{\linewidth}\raggedright
Quantisation into PatternTurn
\end{minipage} & \begin{minipage}[b]{\linewidth}\raggedright
Implemented
\end{minipage} & \begin{minipage}[b]{\linewidth}\raggedright
Implementation chapter
\end{minipage} \\
\begin{minipage}[b]{\linewidth}\raggedright
Rhythmic analysis modules
\end{minipage} & \begin{minipage}[b]{\linewidth}\raggedright
Implemented
\end{minipage} & \begin{minipage}[b]{\linewidth}\raggedright
Core technical contribution
\end{minipage} \\
\begin{minipage}[b]{\linewidth}\raggedright
Response planning
\end{minipage} & \begin{minipage}[b]{\linewidth}\raggedright
Implemented
\end{minipage} & \begin{minipage}[b]{\linewidth}\raggedright
Core technical contribution
\end{minipage} \\
\begin{minipage}[b]{\linewidth}\raggedright
ChucK/Chunity playback
\end{minipage} & \begin{minipage}[b]{\linewidth}\raggedright
Implemented
\end{minipage} & \begin{minipage}[b]{\linewidth}\raggedright
Implementation chapter
\end{minipage} \\
\begin{minipage}[b]{\linewidth}\raggedright
Loop closure after playback
\end{minipage} & \begin{minipage}[b]{\linewidth}\raggedright
Implemented
\end{minipage} & \begin{minipage}[b]{\linewidth}\raggedright
Closed-loop architecture contribution
\end{minipage} \\
\begin{minipage}[b]{\linewidth}\raggedright
Debugging, observability tools, Automated tests and diagnostics
\end{minipage} & \begin{minipage}[b]{\linewidth}\raggedright
Implemented
\end{minipage} & \begin{minipage}[b]{\linewidth}\raggedright
Evaluation/supporting infrastructure
\end{minipage} \\
\begin{minipage}[b]{\linewidth}\raggedright
User evaluation with drummers
\end{minipage} & \begin{minipage}[b]{\linewidth}\raggedright
Deferred
\end{minipage} & \begin{minipage}[b]{\linewidth}\raggedright
Evaluation limitation / future work
\end{minipage} \\
\begin{minipage}[b]{\linewidth}\raggedright
Full response realisation from plan/skeleton
\end{minipage} & \begin{minipage}[b]{\linewidth}\raggedright
Deferred
\end{minipage} & \begin{minipage}[b]{\linewidth}\raggedright
Future work
\end{minipage} \\
\begin{minipage}[b]{\linewidth}\raggedright
Long-term musical memory
\end{minipage} & \begin{minipage}[b]{\linewidth}\raggedright
Deferred
\end{minipage} & \begin{minipage}[b]{\linewidth}\raggedright
Future work
\end{minipage} \\
\midrule\noalign{}
\endhead
\bottomrule\noalign{}
\endlastfoot
\end{longtable}


\section{Contributions}\label{contributions}

A summary of the main contributions is shown in Table~\ref{tab:contributions-summary}.

\begin{longtable}[]{@{}lll@{}}
\caption{Main project contributions}\label{tab:contributions-summary}\\
\toprule\noalign{}
\textbf{Contribution} & \textbf{Realisation (System Components)} &
\textbf{Significance} \\
\textbf{Closed-loop turn-taking interaction model} & TurnLoopController,
TurnPhase, loop closure logic & Formalises real-time human-AI rhythmic
exchange as an iterative, state-driven interaction system \\
\textbf{Real-time sensor-to-symbolic capture pipeline} & Bela input, OSC
messaging, OscHitReceiver, HitBuffer & Enables low-latency
transformation of physical performance into analysable rhythmic data \\
\textbf{Hybrid symbolic-temporal rhythm representation} & TurnWindow,
PatternCompiler, PatternTurn & Captures both quantised structure and
expressive timing, supporting meaningful musical analysis \\
\textbf{Feature-driven rhythmic analysis framework} & Density, energy,
anchor, end-activity, segment profile analysers & Provides
interpretable, multi-dimensional characterisation of rhythmic input \\
\textbf{Deterministic response planning model} & ResponsePlanner,
ResponsePlan, response-type system & Enables explainable mapping from
musical features to structured response strategies \\
\textbf{Structured response generation pipeline} & SkeletonBuilder, map
builders, generation pipeline stages & Decomposes rhythm generation into
controllable stages, enabling variation and constraint-driven
synthesis \\
\textbf{Integrated real-time playback system} & IT4ChuckTurnPlayer,
ChucK/Chunity scripts & Closes the interaction loop through synchronised
audio response \\
\textbf{System observability and debugging framework} & Debug
presenters, visualisation panels & Supports interpretability,
development and evaluation of system behaviour \\
\midrule\noalign{}
\endhead
\bottomrule\noalign{}
\endlastfoot
\end{longtable}
